% Original schematic based on Tutorial 11 Q3 and Lecture 11 pp. 67--70.
\begin{tikzpicture}[
 font=\small,>=Stealth,
 box/.style={draw=noteBlue,fill=noteBlue!4,rounded corners,
 minimum width=1.8cm,minimum height=.8cm,align=center},
 flow/.style={->,thick,noteBlue}
]
\node[box] (x) at (0,0) {图像输入\\784};
\node[box] (h1) at (2.6,0) {$h_1$\\625};
\node[box] (h2) at (5.2,0) {$h_2$\\100};
\node[box,minimum width=2.4cm] (dec) at (8.4,0) {两层解码器\\$100\to625\to784$};
\node[box] (rec) at (11.6,0) {重建图像\\目标为 $x$};
\node[box,draw=ntuRed,fill=ntuRed!4,minimum width=2.4cm] (cls) at (8.4,-1.6) {分类层\\Linear + Softmax};
\node[box,draw=ntuRed,fill=ntuRed!4] (prob) at (11.6,-1.6) {10 类概率\\目标为 $y$};
\draw[flow] (x)--node[above] {$E_1$}(h1);
\draw[flow] (h1)--node[above] {$E_2$}(h2);
\draw[flow] (h2)--(dec);
\draw[flow] (dec)--(rec);
\draw[->,thick,ntuRed] (h2.south)--(5.2,-1.6)--(cls.west);
\draw[->,thick,ntuRed] (cls)--(prob);
\end{tikzpicture}
